Senise töö kõige olulisem metoodiline järeldus on, et äratussõna mudeli arenduses võib pealtnäha hea tulemus olla eksitav, kui hindamismetoodika on puudulik. Lühikeste klippide head tulemused ei taga, et mudel käitub õigesti pideva helivoo režiimis. See probleem ilmnes selgelt hetkel, kui varasemad \enquote{Kratt} tulemused osutusid klipitasemel vaates mõistlikuks, kuid voogedastushindamine ei sisaldanud sisulist taustaheli komponenti.

\section{Avaliku kontrollkatse roll}
Avaliku marvin-kontrollkatse tegemine enne eestikeelse äratussõna juurde liikumist oli teadlik riskide vähendamise samm. Ilma selle vahekontrollita oleks iga hilisem nõrk tulemus võinud näida kui eesti keele, väikese andmestiku või seadme mikrofoni probleem, kuigi tegelik põhjus võinuks peituda hoopis treeninguahela tehnilises piirangus --- täpsemalt samas voogedastushindamise taustaheli komponendi puudujäägis, mille kontrollkatse lõpuks nähtavale tõi.

\section{Kahe raamistiku diagnostiline võrdlus}
\texttt{microWakeWord} valiti põhiraamistikuks ESP32-S3 sihtseadme ja TFLite-väljundi tõttu; \texttt{openWakeWord}i hinnati diagnostilise alternatiivina samadel hindamisperekondadel \cite{microwakeword2026,openwakeword2026}. Tulemuste peatüki tabel~\ref{tab:openwakeword-diagnostic} näitab, et \texttt{openWakeWord} ei andnud selles katses praktilist asenduspunkti: rangemad läved langetasid taustaheli FAPH-i, kuid tegid seda ebaühtlase või kokku variseva tuvastamismäära ning jätkuvalt kõrge segiajamisfraaside FPR-i hinnaga.

Järeldus on metoodiline, mitte raamistikku järjestav. Diagnostiline katse toetab sama tõlgendust, milleni jõudsid \texttt{microWakeWord}i kontrollpunkti- ja konsensuskatsed: suurem või teistsugune mudelirada ei kõrvalda puhaste ning mitmekesiste positiivsete näidete ega fraasiselektiivsete negatiivnäidete nappust. Lisaparameetrid vajavad vastavat andmekatet; vastasel juhul liigub mudel lihtsalt sama FAPH-i, tuvastamismäära ja fraasitäpsuse kompromissi teise punkti.

\section{Domeeninihke mõju}
Kui raamistikuvalik on paigas, nihkub peamine jääkrisk tööriistadelt andmetele --- täpsemalt treening- ja kasutuskeskkonna akustilisele erinevusele. Eestikeelse \enquote{Kratt} mudeli puhul tuleb arvestada, et mudel võib õppida ära salvestusseadme, ruumi või kõneleja eripära, mitte üksnes äratusfraasi enda. Selline domeeninihe on eriti oluline väikese andmestikuga projektis, kus iga allika eripära võib mudeli käitumist ebaproportsionaalselt mõjutada.

\subsection{Andmestiku põhipiirangud}
Käesoleva töö positiivsed näited on kogutud sama tüüpi sihtseadmega, millel mudelit hiljem kasutada soovitakse. See vähendab telefoni ja kaugmikrofoni vahelist domeeninihke riski, kuid ei kõrvalda kõnelejate mitmekesisuse probleemi. Varajane käsitsi kogutud andmestik põhineb kitsal kõnelejate valimil ning sünteetilise kõnega laiendamine ei asenda täielikult tegelikke kasutajaid. Sünteetiline kõne aitab katta hääletämbreid ja kõnetempot, kuid võib samal ajal anda liiga puhta või liiga ühtlase hinnangu mudeli tuvastamismäärale.

Negatiivse klassi puhul on eraldi tähtis eristada kahte riski. Üks risk on üldine taustaheli, kus äratusfraasi ei öelda. Teine risk on foneetiliselt sarnased fraasid, mis kõlavad äratussõnale lähedaselt, kuid ei peaks mudelit käivitama. Kui need kaks hinnatakse üheainsa mõõdikuna kokku, võib mudel paista tugev, kuigi ta ei erista fraasi \enquote{Kuule Kratt} piisavalt hästi sarnastest ütlustest.

\subsection{Hindamise ja praktilise kasutuse põhimõtted}
Treeningu detailsete parameetrite asemel on lõputöö seisukohalt olulisem põhimõte, et igat mudelit tuleb hinnata mitmes eri olukorras. Klipitaseme täpsus näitab, kas mudel tunneb sihtfraasi isoleeritud näites ära. Taustaheli FAPH näitab, kui sageli mudel valesti käivitub pidevas kuulamisrežiimis. Sarnaste fraaside test näitab, kas mudel õppis ära kogu fraasi või ainult selle alguse. Reaalajas seadmel testimine näitab, kas laboritingimustes valitud lävi säilib tegeliku mikrofoni, kauguse ja ruumi korral.

Konsensushindamine, kus äratus otsustatakse kahe mudeli ühise signaali põhjal, osutus üheks viisiks vähendada valeaktiveeringuid. Selle tulemuse tõlgendamisel tuleb siiski rõhutada, et madal valeaktiveeringute määr ühel taustaheli rajal ei ole iseseisvalt piisav alus praktiliseks kasutuseks. Kui sama seadistus kaotab tegelike kõnelejate tuvastamismäära või jääb tundlikuks sarnastele fraasidele, ei ole mudel praktiliselt parem.

Seadmel käitamise seisukohalt kinnitab töö, et mudelit saab käivitada väikesel mikrokontrolleripõhisel häälseadmel ning siduda Home Assistanti hääljuhtimise ahelaga. See on oluline praktiline tulemus, kuid arutelu peatükis ei ole keskne seadistuse täpne kirjeldus. Pärast äratussõna tehnilise toimivuse valideerimist jääb järgmine süsteemitaseme küsimus kasutusolukorra kontrolliks: kas äratus käivitub õigel fraasil, kui kasutaja räägib eri kauguselt, ning kas valeaktiveeringuid tekib piisavalt harva, et laiem häälkäskude ahel oleks päriselt kasutatav.


\section{Standardsete võrdlusaluste ebapiisavus reaalse kasutuse ennustamisel}
\label{sec:benchmark-gap}

Käesoleva töö üks kõige selgemaid ja praktiliselt olulisemaid leide on korduvalt täheldatud lahknevus standardsete KWS-võrdlusaluste tulemuste ja tegeliku kasutajakogemuse vahel.

\subsection{Empiiriline tõendus lahknevusest}

Kõige teravamalt ilmnes probleem kahe mudeliversiooni võrdluses. Üks versioon paistis standardsete võrdlusaluste järgi kasutuselevõtuks sobiva mudelina: sellel oli madal FAPH, nullilähedane klipitaseme tõrkemäär ja kõrge tuvastamismäär sünteetilistel positiivsetel klippidel (vt ptk~\ref{chapter:results}, tabel~\ref{tab:checkpoint-headline}). Reaalajas testimine näitas aga vastupidist pilti: mudel aktiveerus sageli sarnastel fraasidel nagu \enquote{kuule rott} ja \enquote{kuule kraam} ning tegeliku kõneleja häälel jäi tuvastamismäär oodatust madalamaks.

Seevastu teine versioon sai standardsetel võrdlusalustel selgelt nõrgema tulemuse, kuid käitus reaalses proovikohas mõnes stsenaariumis paremini: päris kõneleja häälel käivitus see sagedamini õigel ajal, sarnastel fraasidel aktiveerus harvemini ning mittekõnelised helid tekitasid vähem valekäivitusi. Täpsemad arvud on esitatud tulemuste peatükis; arutelu seisukohalt on oluline just suunaerinevus, mitte konkreetne mudelitähis.

Sama suunaerinevus kordus mitme vaheversiooni puhul kitsa reaalkõnelejate baasi peal: parem võrdlusalustulemus ei taganud paremat päriskõneleja tuvastamismäära. Seetõttu on tegu suunanäitajaga, mitte kinnitatud mustriga, ja laiema valimiga kasutajatest jääb vajalikuks täienduseks.

\subsection{Klipitaseme tuvastamismäär ja reaalne kõneleja}

Positiivse tuvastamismäära mõõtmine toimus XTTS-kloonitud hääle isoleeritud klippide peal. Kõik klipid olid normaliseeritud, konstantse kauguse ja helitugevusega ning samalt kõnelejalt kloonitud. Selline testimine andis väljundväärtuseks täpselt 1,0000, mis tekitas vale kindlustunnet. Tegelikud kõnelejad varieeruvad häälduse, helitugevuse ja kauguse poolest märkimisväärselt. Reaalsete ütluste silutud väljundid jäid katseprotokolli alusel valitud otsuseläve lähedale; täpne ütluse-tasandi jaotus on edaspidi kasutajatesti tulemuste osa.

Park et al.~\cite{park2024adversarial} on näidanud, et TTS-andmetega treenitud mudelid kipuvad sünteesitud kõne peal näitama ülehinnatud tuvastamismäära, sest TTS-kõne ei sisalda sama varieeruvust kui loomulik kõne. Käesoleva töö tulemused kinnitavad seda leidu, kuid lisavad olulise täpsustuse: probleem ei piirdu ainult treenimisega, vaid laieneb ka hindamisele. Kui testikomplekt koosneb samadest TTS-häältest, millega treeniti, on tuvastamismäära hinnang topelt üle paisutatud.

\subsection{Tihe kõnekorpus versus reaalne keskkond}

FAPH mõõtmine Common~Voice eestikeelsel kõnel kujutab olukorda, kus mikrofon kuuleb pidevat, selgesti artikuleeritud kõnet. Tegelikus nutikodu kasutuses on heli koostis aga hoopis teine: valdavalt vaikus, vahetevahel televisioon, muusika, vestlused teises toas, köögihelid. Mudel, mis tihedal kõnel aktiveerub 243 korda tunnis, ei pruugi elutoas aktiveeruda üldse, sest suurem osa ajast puudub foneetiliselt sarnane sisend. Vastupidi, mudel, mis tihedal kõnel näitab FAPH\,=\,14,40, võib olla liialt tundlik just nende üksikute mitte-kõne mustrite suhtes, mida standardne testkorpus ei sisalda.

Dubois et al.~\cite{dubois2020triggers} demonstreerisid seda lahknevust empiiriliselt: kommertsiaalsetel nutikõlaritel mõõdeti kuni 0,95 tahtmatut aktiveerimist tunnis, kui seadmed kuulsid televisioonisignaali --- tulemus, mida standardsed võrdlusalused ei olnud ennustanud. Sch{\"o}nherr et al.~\cite{schoenherr2022accidental} kinnitavad seda mustrit eraldi katsega, näidates, et juhuslike aktiveerumiste määr varieerub seadme ja korpuse vahel märkimisväärselt. Sensory, kes arendab Samsungi äratussõna mudeleid, on tunnistanud, et nende sisemine hindamine ei kattunud alati klientide tegeliku kogemusega~\cite{sensory2024realworld}.

\subsection{Klipi- ja koduheli testide kontekstipiirid}

Foneetiliselt sarnaste fraaside eristamist testiti isoleeritud klippidena. Mudel, mis korrektselt tõrjub üksiku klipi \enquote{kuule rott}, võib reageerida teisiti, kui sama fraas tuleb lause kontekstis: \enquote{kuule, rott on köögis}. Voogedastusrežiimis on mudeli sisemine olek eelmiste kaadrite mõjutatud, koartikulatsioon muudab akustilist signaali ja ajastus varieerub. Klipitaseme hindamine ei kajasta neid efekte.

Teine katvusprobleem puudutab mittekõnelisi koduhelisid. Standardsed KWS-võrdlusalused sisaldavad kõnet ja vahel ka keskkonnahelisid (nt MUSAN-i müra \cite{musan2015}), kuid ei sisalda konkreetseid koduse keskkonna stiimuleid, nagu nõude kolin, kaugelt kostev televisioon või juhuslik muusika. Ükski kasutatud standardne võrdlusalus neid juhtumeid süstemaatiliselt ei kata; just see lõhe motiveeris käesoleva töö stsenaariumpõhist hindamisvoogu, mida kirjeldab järgmine alajaotus.

\subsection{Stsenaariumpõhine testimisvoog ja Androidi väliterminal}
\label{sec:android-logger}

Eelnev analüüs viitab sellele, et äratussõnamudelite usaldusväärne hindamine nõuab standardsete mõõdikute kõrvale täiendavat \emph{stsenaariumpõhist testimisvoogu}: pidevat reaalse kasutuse helivoogu, milles äratussõna esinemist on võimalik kontrollida. Pika annoteeritud koduheli kogumine, salvestamine ja säilitamine on aga selles töös kahel praktilisel põhjusel keeruline: mahukate WAV-failide salvestus tekitab failimahu ja varundamise koormuse, ning pideva mikrofoni-streami salvestamine kogu autori päeva jooksul tõstab privaatsusriski, mida ei lahenda ei kohalik krüpteerimine ega lühike säilitusaken üksi.

Käesolevas töös ehitati selle takistuse kõrvale hiilimiseks Androidi väliterminal-rakendus, mis jooksutab \emph{paralleelselt mitut äratussõnamudelit} reaalse mikrofoni-streami peal kogu päeva vältel, ilma helisalvestust säilitamata. Salvestatakse ainult mudeli aktiveerumise sündmused (mudeli identifikaator, ajatempel, otsustusskoor) sellistel ajaakendel, kus sihtfraasi \enquote{Kuule Kratt} \emph{ei} ole öeldud --- need on definitsiooni järgi valeaktiveeringud. Privaatsusprobleem laheneb seetõttu rakenduse disainiloogika, mitte hilisema redaktsiooni kaudu: helivoog ei lahku seadme muutmälust ja ükski heliklipp ei jää püsivasse mällu. Sama lahendus kaotab samuti annotatsioonikoormuse, sest sündmuste \enquote{tõene}/\enquote{vale}-silt sõltub ainult sellest, kas autor ütles selles aknas sihtfraasi --- ja autor ei ütle seda päeva jooksul tavakasutuses.

Logija praktiline panus oli üks selle töö kõige selgemini mõõdetavaid tagasiside-tsükleid: kogutud valeaktiveeringud (lauapealsetel löökidel, klaviatuuriklobinal ja teatud kõnemustritel) lisati järgmise treeningu negatiivklassi, ning sama klassi valeaktiveeringud kadusid järgmise mudeliversiooni FAPH-mõõdikutest praktiliselt täielikult. Sellise tsükli formaalne mõju on selge: pidevkogumise põhine negatiivide laiendamine annab FAPH-mõõdikus piisavalt suuri parandusi, et tasub eraldi metoodilise sammuna registreerida. Mehhanism ise --- pikaajaline välipõhine valeaktiveeringute kogumine reaalselt seadmelt --- ei ole äratussõna kirjanduses uudne \cite{apple-heysiri2017,sigtia2020multitask}, mistõttu käesolev töö ei väida selle leiutamist, vaid esitab Androidi väliterminalit konkreetse ja sellele tööle kohandatud teostusena, mille kogumislogi sai mudeliarendusele otsene sisend.

Lahendus toob endaga ühe metoodilise kompromissi: kuna salvestust ei säilitata, ei ole varasema väliterminali-mõõdiku taasmäng uuemate mudeliversioonidega võimalik --- iga uus mudel nõuab uut paralleelset jooksu samadel käitumistingimustel, et FAPH-numbreid saaks võrrelda. Selle hinnaga tuleb arvestada igas pikaajalises pidevkogumise põhises hindamiskihis ning see motiveerib multimodaalsete annoteeritud kodukõnekorpuste \cite{chen2022misp} kasutamist sealsamas valideerimisahelas, kus reprodutseeritav taasmäng on prioriteet.

\section{Agentpõhine arendus töövõimendajana}

Käesoleva töö üks keskseid järeldusi on, et tänapäevane agentpõhine tööviis nihutab äratussõna arenduse peamise pudelikaela rakenduselt tõendusmaterjalile: üksi töötav autor saab ehitada lühikese ajaga oluliselt laiema tehnilise süsteemi kui varem, kuid lõppjärelduste tugevus sõltub endiselt pärisandmete kogumisest, hindamismetoodikast ning esialgsete tulemuste eristamisest lõplikest väidetest.

Agentpõhine arendus võimaldas selles töös kolm tegevust, mis oleksid üksiku autori jaoks muidu jäänud mahult kättesaamatuks: Androidi-põhine valeaktiveeringute logija päris kasutusolukorra negatiivnäidete kogumiseks, stsenaariumipõhine hindamisahel sõltumatu taustaheli FAPH mõõtmiseks ning paralleelselt arendatud hindamis- ja analüüsiskriptid, mis lubasid avastada andmelekkeid ja korrata hindamist sõltumatutel testikomplektidel sama nädala jooksul. Need ei ole pelgalt tootlikkuse näitajad, vaid uurimist sisuliselt kujundavad tegurid: ilma logijata poleks taustaheli FAPH-i saanud kalibreerida päris koduakustika peal ning ilma kiire iteratsioonita poleks andmelekked jõudnud avastatuks.

Mõju ei tohi siiski üle tõlgendada. Tehisagendid vähendavad rakendusliku töö, skriptimise ja prototüüpimise kulu, kuid ei vähenda samal määral tõendusmaterjali hankimise kulu: nad ei loo juurde päriskõnelejaid, ei asenda sõltumatuid testikomplekte ega lahenda välise valiidsuse probleemi. Just see asümmeetria --- süsteemiehituse kiirenemine ilma tõendusmaterjali kiirenemiseta --- on käesoleva peatüki teesi tehniline alus.

\section{Edasised suunad}
\label{sec:future-directions}

Käesolev töö avab kolm konkreetset jätkusuunda. Esimene puudutab arhitektuuri --- kuidas laiendada üheastmelist äratussõnamudelit järelkontrolli sisaldavaks kaskaadiks. Teine puudutab metoodika rangust --- kuidas vähendada ühe treeningujooksu ja kitsa kõnelejavalimi loodud määramatust korduskatsete ja kõneleja-lahuse jaotuse kaudu. Kolmas puudutab valideerimist --- millised mõõtmised tõstaksid usaldust tulemuste vastu rohkem kui täiendavad mudeliversioonid.

\subsection{Kaskaadarhitektuur}
\label{sec:future-cascade}

Konsensus, mida käesolev töö dokumenteerib --- kaks ESP32-S3-klassi mudelit, mille väljundid ühitatakse loogilise konjunktsiooniga (§\ref{sec:expert-consensus}) ---, on erijuhtum laiemast \emph{kaskaadarhitektuurist}, kus väike alati-aktiivne detektor pakub esialgse kandidaadi ja teine, ressursimahukam aste kontrollib seda põhjalikumalt. Kaskaadi idee on äratussõna kirjanduses levinud \cite{gruenstein2017cascade,apple_voice_trigger_2023}; seda on rakendatud nii üldistes mobiilsetes KWS-süsteemides \cite{michaely2017googlekws,sigtia2020multitask} kui ka madala latentsi voogedastuses \cite{garg2021streaming}.

Krati süsteemis on kaskaadi riistvarapaigutus loomulik. Esimene aste mahub juba praegu ESP32-S3 mikrokontrolleri mälupiiridesse: \texttt{v16c} on $148$\,KB kvantiseeritud TFLite-mudel pluss ${\sim}50$\,KB \texttt{tensor\_arena} tööpuhver (§\ref{subsec:quantization}). Teise astme saaks asetada Home Assistanti hostarvutile, kus on rohkem arvutusressursse ja vähem reaalajapiirangut, mistõttu see ei pea olema sama kitsa eelarvega kui esimene. Teise astme funktsionaalne roll on praktikas kitsalt määratletav: kontrollida, kas mikrokontrolleri aktiveerus tuleneb tõepoolest täisfraasist \enquote{Kuule Kratt} või \enquote{Kule Kratt}, mitte üksiku eesliite \enquote{kuule}/\enquote{kule} hääldusest. Need on täpselt veakategooriad, mida v17 mudelid demonstreerisid (§\ref{sec:positive-quality-control}) ja mida fraasistruktuuri kontrollivad mõõdikud (§\ref{sec:phrase-selectivity-metrics}) eraldi mõõdavad.

Suuna rakendatavus sõltub kahest tingimusest. \emph{Esiteks} peab teine aste opereerima esimesest erineval esitusel --- näiteks foneemitasemel tunnustel või lühikese järel-ASR-päringuga \cite{wang2020lfmmi,rikhye2021personalized} ---, vastasel juhul õpib see samad lühiteed, mis on käesolevas töös ka esimese astme juures dokumenteeritud (vt §\ref{sec:metric-iteration-conclusion}), ega lisa tegelikku eristusvõimet. \emph{Teiseks}: kui esimese astme tuvastamismäär reaalsetel \enquote{Kule}-hääldustel on juba madal (§\ref{sec:checkpoint-listening}), kustutab järelkontroll tõelisi äratusi, mitte valeaktiveeringuid \cite{kumar2020wakeword}; järelkontrolli põhjendab vaid see, kui esimene aste säilitab kõrge tuvastamismäära ja aktiveerub liiga tihti. Käesoleva töö lõppolukord --- $0{,}79$ valeaktiveeringut tunnis 3,82\,h pikkusel Common~Voice ET voogedastusrajal (§\ref{sec:expert-consensus}) --- annab kaskaadi sisendiks juba madala FAPH-i; edasise suuna empiiriline kasu sõltub seetõttu vähem FAPH-i edasisest langetamisest ja rohkem sellest, kas teine aste suudab tagasi võita reaalsete kõnelejate tuvastamismäära ilma esimese astme künnist liiga madalale lükkamata.

\subsection{Metoodika edasiarendus}
\label{sec:future-methodology}

Käesoleva töö ablatsioonid (residuaalühendused §\ref{sec:residual-ablation}, SpecAugment §\ref{sec:specaug-ablation}) põhinevad ühel treeninguseemnel, mistõttu tähelepanekud kannavad mudeli initsialiseerimise ja partii-järjekorra varieeruvust, mida üks jooks ei eralda mõju enda suurusest. Kõige otsesem rangusparandus on iga raporteeritud ablatsiooni kordamine vähemalt kolme--viie eri RNG-seemnega ja efekti suuruse esitamine koos seemnete-vahelise standardvea või bootstrap-vahemikega; sama põhimõte laieneb kontrollpunkti-FAPH eksperimentidele (§\ref{sec:checkpoint-objective-control}), kus praegu pole võimalik eristada \texttt{target\_minimization} valikut konkreetse jooksu juhuslikkusest.

Positiivse klassi kõnelejavalimi laiendamine on selle metoodika teine pool. Käesoleva töö üldistuvuse väiteid piirab eelkõige üks treeningust eraldi hoitud reaalne kõneleja (Kõneleja~D, $N=145$); juba 5--10 sõltumatu täiskasvanud kõneleja kaasamine kõneleja-disjoint testikomplektina annaks Wilsoni 95\%~UV-d, mis lubavad eristada üksiku punktivõrdluse asemel mudelivahelisi erinevusi populatsioonitasandil. Kaasamise loogiline järgmine samm on lisada kõneleja vanus, sugu ja aktsent registreeritud metaandmetena, et tuvastamismäära raporteerimine saaks edaspidi toimuda alarühmade kaupa.

\subsection{Hindamise ja valideerimise täiendused}
\label{sec:future-validation}

Esimene konkreetne valideerimissamm on külmutatud lävedega kasutajapiloot reaalses Korvo-2-seadmes (vt §\ref{sec:user-test-methodology}): praegu raporteeritud mõõdikud põhinevad skriptitud taasmängul, mistõttu süsteemitaseme käitumist --- latentsus, mikrofonigeomeetria mõju, kasutaja ootuste lõhe --- ei saa neist otse järeldada. Teine samm on praeguse \texttt{faph\_cv\_et} kõrvalejäetud rea ümberkonstrueerimine \emph{kõneleja-disjoint} jaotusega: Common~Voice eestikeelses jaotuses tuleb tagada, et ükski kõneleja ei esine korraga treeningus ja kõrvalejäetud komplektis, vastasel juhul mõõdab voogedastus-FAPH osaliselt sama kõneleja-mikrofoni allkirja tagasi.

Kolmas samm on tehtud \texttt{openWakeWord}i diagnostilise läbiproovi laiendamine täielikuks raamistike võrdluseks: mõlemas raamistikus tuleks treenida sama andmestiku ja eelregistreeritud lävepoliitikaga mudelid ning hinnata neid samadel kõrvalejäetud komplektidel. Käesolev läbiproov on piisav järelduseks, et raamistiku vahetamine üksi ei eemalda andmekatte ja fraasiselektiivsuse puudust, kuid mitte piisav üldiseks \texttt{microWakeWord}i ja \texttt{openWakeWord}i paremusjärjestuseks. Neljas samm on lõpliku tegelikult eraldatud (\emph{true held-out}) testikomplekti määratlemine, mida ei kasutata mudelivaliku ega lävekülmutamise ühegi otsuse juures, vaid raporteeritakse alles töö lõpus üks kord; käesoleva töö \texttt{faph\_cv\_et} on lävede ja konsensuspaaride iteratsioonide tõttu muutunud regressioonikomplektiks ning ei sobi enam selle rolli täitma.
